% !TEX root = ../main.tex
% ============================================================================
% Budget: 1.5 pages, with Figure 2 (pipeline) and Table 1 (inputs).
%
% This section says what the method IS, in the present tense, so that a reader
% could build it again. It holds no results, no cluster detail, and no history
% of how the team got here: results go to Section 5, the cache and the cluster
% to Appendix C, the differences from the proposal to Appendix A.
%
% The prose of 3.2 to 3.4 is carried over from the midterm report (Method as
% Implemented). The equations of 3.1 are new and were written from
% confhead/src/confhead/labels.py, geometry.py (Fit, fit_scale_shift) and
% head.py (masked_bce); the owners should check them against the code.
% ============================================================================
\section{Method}
\label{sec:method}

Figure~\ref{fig:pipeline} shows the pipeline from the raw data to the loss. Only the head is trained; everything to its left is computed once per image and cached.

\begin{figure*}[t]
  \centering
  \resizebox{\textwidth}{!}{% Pipeline of the confidence head, condensed from Docs/Wiki/Models/run1/run1.canvas.
% Needs tikz with the libraries positioning and arrows.meta (loaded in main.tex).
\begin{tikzpicture}[
    font=\scriptsize,
    box/.style={draw=black!45, rounded corners=1.5pt, align=center, text width=2.55cm,
                minimum height=0.82cm, inner sep=2pt},
    raw/.style={box, fill=orange!18},
    calc/.style={box, fill=violet!12},
    cache/.style={box, fill=cyan!15},
    chan/.style={box, fill=green!18},
    lab/.style={box, fill=yellow!30},
    head/.style={box, fill=black!6},
    arr/.style={-{Stealth[length=1.6mm]}, draw=black!60, line width=0.4pt},
    colhead/.style={font=\scriptsize\bfseries, align=center, text width=2.9cm},
  ]
  % Columns: x = 0, 3.25, 6.5, 9.75, 13. Rows: y = 0, -1.05, -2.1, -3.15, -4.45.
  \node[colhead] at (0, 0.85)     {nuScenes data};
  \node[colhead] at (3.25, 0.85)  {Cache build\\(once per image)};
  \node[colhead] at (6.5, 0.85)   {Cache entry};
  \node[colhead] at (9.75, 0.85)  {Head input\\(set per run)};
  \node[colhead] at (13, 0.85)    {Trained};

  % Raw data
  \node[raw] (rimg) at (0, -0.525) {Camera keyframe\\$1600 \times 900$, 6 cameras};
  \node[raw] (rnbr) at (0, -2.1)   {Neighbor frames\\$\pm 1, \pm 2$ at 12\,Hz};
  \node[raw] (rrad) at (0, -3.15)  {Radar points\\5 sweeps, no elevation};
  \node[raw] (rlid) at (0, -4.45)  {Lidar points\\never an input};

  % Cache build
  \node[calc, minimum height=1.87cm] (da2) at (3.25, -0.525) {Frozen Depth\\Anything V2 ViT-B\\at $924 \times 518$};
  \node[calc] (btmp) at (3.25, -2.1)  {4 extra passes; follow each cell's 3D point};
  \node[calc] (brad) at (3.25, -3.15) {Project to camera;\\radar fit};
  \node[calc] (blid) at (3.25, -4.45) {Project to camera;\\lidar fit};

  % Cache entry
  \node[cache] (cfeat) at (6.5, 0)     {Decoder features\\\texttt{output\_conv1}, fp16};
  \node[cache] (cdisp) at (6.5, -1.05) {Disparity grid};
  \node[cache] (ctmp)  at (6.5, -2.1)  {Temporal spread};
  \node[cache] (crad)  at (6.5, -3.15) {Radar points, fit\\(scale, shift)};
  \node[cache] (clid)  at (6.5, -4.45) {Lidar points, fit\\(scale, shift)};

  % Head input and labels
  \node[chan] (ifeat) at (9.75, 0)     {Features (64)};
  \node[chan] (idep)  at (9.75, -1.05) {Depth (1)\\disparity / image mean};
  \node[chan] (itmp)  at (9.75, -2.1)  {Temporal spread (1)};
  \node[chan] (irad)  at (9.75, -3.15) {Radar residual and validity (2)};
  \node[lab]  (ilab)  at (9.75, -4.45) {Target $\mathbf{1}[\mathrm{AbsRel} > \tau]$ and label mask};

  % Head and loss
  \node[head, minimum height=3.97cm] (head) at (13, -1.575)
    {\textbf{Confidence head}\\[2pt]standardize channels\\$3 \times 3$ conv, 64\\$3 \times 3$ conv, 64\\$3 \times 3$ conv, 32\\$1 \times 1$ conv\\[2pt]$p_\mathrm{error}$ per grid cell};
  \node[lab] (loss) at (13, -4.45) {Masked BCE on labeled cells};

  % Edges
  \draw[arr] (rimg) -- (da2);
  \draw[arr] (da2.east |- cfeat) -- (cfeat);
  \draw[arr] (da2.east |- cdisp) -- (cdisp);
  \draw[arr] (rnbr) -- (btmp);
  \draw[arr] (rrad) -- (brad);
  \draw[arr] (rlid) -- (blid);
  \draw[arr] (btmp) -- (ctmp);
  \draw[arr] (brad) -- (crad);
  \draw[arr] (blid) -- (clid);
  \draw[arr, dashed] (blid.west) -- ++(-0.22, 0) |- (btmp.west);
  \draw[arr] (cfeat) -- (ifeat);
  \draw[arr] (cdisp) -- (idep);
  \draw[arr] (ctmp) -- (itmp);
  \draw[arr] (crad) -- (irad);
  \draw[arr] (clid) -- (ilab);
  \draw[arr] (ifeat) -- (ifeat -| head.west);
  \draw[arr] (idep) -- (idep -| head.west);
  \draw[arr] (itmp) -- (itmp -| head.west);
  \draw[arr] (irad) -- (irad -| head.west);
  \draw[arr] (ilab) -- (loss);
  \draw[arr] (head) -- node[right, font=\tiny] {logits} (loss);
\end{tikzpicture}
}
  \caption{The pipeline. Orange: raw nuScenes data. Violet: a computation done once per image when the cache is built. Cyan: an array stored in the cache. Green: an input group of the head, which a run can switch off (Table~\ref{tab:input-ablation}). Yellow: the lidar labels and the loss. The dashed arrow marks that the temporal spread uses the lidar fit to lift each cell to a 3D point.}
  \label{fig:pipeline}
\end{figure*}

% ----------------------------------------------------------------------------
\subsection{Problem Formulation}
\label{sec:formulation}
% New for the final report (checklist B5: numbered equations).
% \todo for the owners: check each equation against the code named above.

The depth model returns a disparity map $q$ that is inverse depth up to an unknown scale and shift. For an image with lidar returns $i = 1, \dots, N$ at depths $z_i$, let $q_i$ be the disparity sampled at the pixel of return $i$. We fit the two unknowns per image by least squares against inverse depth,
\begin{equation}
  (s, t) = \operatorname*{arg\,min}_{s,\,t} \sum_{i=1}^{N} \Bigl( s\,q_i + t - \frac{1}{z_i} \Bigr)^{2},
  \label{eq:fit}
\end{equation}
and call the fit failed when $s \le 0$ or the returns are too few. The predicted depth at a return clips the fitted inverse depth to the range of the labels before inverting it,
\begin{equation}
  \hat d_i = \Bigl[ \operatorname{clip}\bigl( s\,q_i + t,\ 1/z_{\max},\ 1/z_{\min} \bigr) \Bigr]^{-1},
  \label{eq:depth}
\end{equation}
with $z_{\min} = 1$\,m and $z_{\max} = 80$\,m. The error of the depth at a return is its absolute relative error,
\begin{equation}
  e_i = \frac{|\hat d_i - z_i|}{z_i}.
  \label{eq:absrel}
\end{equation}
Because of Eq.~\ref{eq:fit}, $e_i$ is the error under the best scale and shift of that image: it measures the shape of the depth, and a frame whose overall scale is off still counts as correct.

Let $c$ index the cells of the output grid, let $M$ be the set of cells that receive at least one return, and let $e_c$ be the error of the nearest return in cell $c \in M$. The label is
\begin{equation}
  y_c = \mathbf{1}\,[\, e_c > \tau \,], \qquad \tau = 0.10 .
  \label{eq:label}
\end{equation}
The head $f_\theta$ maps the input $x$ (Section~\ref{sec:inputs}) to one logit per cell, and
\begin{equation}
  p_c = \sigma\bigl( f_\theta(x)_c \bigr)
  \label{eq:perror}
\end{equation}
is its estimate of the probability that $y_c = 1$. The confidence of cell $c$ is $1 - p_c$. The head is trained on the cells in $M$ and predicts for every cell.

\todo{Check Eqs.~\ref{eq:fit} to \ref{eq:perror} against \texttt{labels.py}, \texttt{geometry.py}, and \texttt{head.py}; state the minimum number of returns for a fit.}

% ----------------------------------------------------------------------------
\subsection{Frozen Backbone and Head Inputs}
\label{sec:inputs}

\paragraph{Backbone.} The backbone is the stock relative Depth Anything V2 ViT-B checkpoint \citep{yang2024depthanythingv2}, frozen and run once per image at $924 \times 518$. Its output is disparity with unknown scale and shift, so each cached entry carries two per-image fits of the form of Eq.~\ref{eq:fit}: one against the keyframe lidar, used for the labels and the temporal channel, and one against the radar points, used for the radar residual.

\paragraph{Inputs.} Everything lives on a $132 \times 74$ grid, the ViT patch grid times two, where one cell covers about $12 \times 12$ pixels of the $1600 \times 900$ camera image. We use times two to fit the storage (about 290\,GB, against 1.1\,TB at times four). The head input has 68 channels (Table~\ref{tab:inputs}), each standardized by its training-set mean and standard deviation. Channel groups can be switched off at load time, which is how the ablation runs of Section~\ref{sec:input-ablation} are made.

\begin{table}[t]
  \centering
  \small
  \setlength{\tabcolsep}{4pt}
  \begin{tabular}{@{}l l p{3.2cm}@{}}
    \toprule
    \textbf{Ch.} & \textbf{Input} & \textbf{How it is made} \\
    \midrule
    0--63 & Decoder features & \texttt{output\_conv1} tap, 64 channels, area-averaged to the grid \\
    64 & Norm.\ disparity & Disparity divided by its image mean; no fit \\
    65 & Temporal spread & Std divided by mean ($\sigma/\mu$) of inverse depth over the keyframe and its neighbors at $\pm1, \pm2$ (12\,Hz; up to 5 frames), ego-motion compensated \\
    66 & Radar residual & $|\hat d - z| / z$ under the radar fit, 5 sweeps, radius 1 cell \\
    67 & Radar validity & 1 where the residual holds a measurement \\
    \bottomrule
  \end{tabular}
  \caption{The 68 input channels of the confidence head. Channel 64 is called the depth channel in the rest of the report. In channel 66, $\hat d$ is the depth that the radar fit gives the model's disparity at a radar point, $z$ is the range of that point, and the nearest point wins a cell where two footprints overlap.}
  \label{tab:inputs}
\end{table}

% ----------------------------------------------------------------------------
\subsection{Lidar-Derived Labels}
\label{sec:labels}

For each lidar return within 1--80\,m, the absolute relative error (AbsRel, Eq.~\ref{eq:absrel}) of the fitted prediction at that pixel is placed in its grid cell; the nearest point wins a cell. Here the prediction is the model disparity sampled at the return's own pixel and passed through the lidar fit (Eq.~\ref{eq:depth}). The target is Eq.~\ref{eq:label}, and a mask marks the labeled cells. About 25\% of cells on the $132 \times 74$ grid carry a label, about 2,450 per image. The lidar returns behind them, about 3,600 per image, cover 0.25\% of pixels. \todo{Refresh the three counts on the complete cache, and say how the label masks were settled: the ego body in CAM\_BACK (D8), bleed-through returns (D9), and images with no lidar fit (D13). These change the labels, so they must be settled before the final scoring (checklist B1).}

% ----------------------------------------------------------------------------
\subsection{Confidence Head and Loss}
\label{sec:head}

The head is three $3 \times 3$ convolutions (stride 1 and padding 1, so the grid size is kept) with ReLU and widths 64, 64, 32, then a $1 \times 1$ convolution to one logit per cell. With all 68 inputs it has 94,657 weights, and its receptive field is $7 \times 7$ cells. The loss is binary cross-entropy (BCE) over the labeled cells only, with positives weighted by the negative-to-positive ratio $w_+$ of the training labels:
\begin{equation}
  \mathcal{L} = -\frac{1}{|M|} \sum_{c \in M} \Bigl[ w_+\, y_c \log p_c + (1 - y_c) \log (1 - p_c) \Bigr].
  \label{eq:loss}
\end{equation}
This weight raises every logit by $\ln w_+$ (\todo{Refresh; 0.506 for $w_+ = 1.658$ on the partial cache; give the value of the complete cache}). The ranking metrics are unaffected, but the calibration shifts, so we report the expected calibration error (ECE) after subtracting it (Section~\ref{sec:protocol}).

% ----------------------------------------------------------------------------
\subsection{Sensors Needed at Inference}
\label{sec:sensors}
% New subsection. It answers the critical item of the midterm review (A1, A2,
% B1 "lidar in the inputs") in one place, so that the results can refer to it.
% Rubric gate 4 checks this subsection.

Training needs lidar for the labels. Whether inference needs it depends on the inputs a head reads. The decoder features and the normalized disparity need no fit. The radar residual uses the radar fit, which a production car can compute. The temporal spread lifts each cell to a 3D point with the lidar fit, so every head that reads it uses lidar at inference and is an upper bound for a deployed system. Two heads are therefore free of lidar: \texttt{m1-features-only}, which reads the decoder features alone and needs no sensor fit at all, and \texttt{m1-no-temporal}. Table~\ref{tab:main} marks each method accordingly.

\todo{State the team's decision on checklist item B1. Either (a) the temporal channel was rebuilt without lidar, or (b) it was kept and the lidar-free heads are presented as the deployable models, with the 68-channel head and \texttt{m2-radar-lidar-fit} as upper bounds. This paragraph assumes (b).}
